%! Characteristics of Logarithmic Functions — Unit 7, Lesson 7.0 — Lesson Plan
\documentclass[10pt]{article}
\usepackage{algebra2-article}
\usepackage{algebra2-boxes}
\usepackage{graphicx}   % -article does NOT load graphicx; needed for thumbnails/screenshots
\graphicspath{{images/}}
\newcommand{\TallMath}[1]{$\displaystyle #1\rule[-1.4em]{0pt}{3.2em}$}
\newcommand{\CourseName}{Algebra 2}
\newcommand{\MeetingLength}{60 minutes}
\newcommand{\UnitNumberName}{Unit 7: Logarithmic Functions \quad}
\newcommand{\LessonNumberName}{Lesson 7.0: Characteristics of Logarithmic Functions}

\begin{document}
\begin{center}
    {\Large\bfseries \CourseName \\
    {\small\bfseries \UnitNumberName \LessonNumberName}}
\end{center}

\begin{tcolorbox}[colback=forestbg, colframe=forest, boxrule=0.9pt, arc=2mm,
    left=2mm, right=2mm, top=1.5mm, bottom=1.5mm]
    \textbf{Primary Objective:} Students can read the graph of a logarithmic function --- both cases,
    $b>1$ and $0<b<1$, taught side by side --- and name every characteristic: the \emph{domain},
    which must be computed by solving \emph{argument} $>0$; the \emph{vertical asymptote} that stands
    at the edge of that domain; and the \emph{range}, \emph{intercepts},
    \emph{increasing/decreasing} and \emph{positive/negative} intervals, absence of \emph{extrema},
    and \emph{end behavior} --- including the honest answer that the left ``end'' is not
    $x\to-\infty$ but $x\to 0^{+}$. Students also explain where the curve came from: $y=\log_b x$ is
    $y=b^x$ reflected over $y=x$, so the two \emph{swap} domain and range, and the exponential's
    horizontal asymptote stands up into the logarithm's vertical one.
    \\[2pt]
    \textbf{Standards (2023 VA SOL):} \textbf{A2.F.2} --- investigate and analyze the characteristics
    of logarithmic functions graphically and algebraically. This Lesson~0 \emph{introduces} the one
    row the \S3 spine marks new for Unit~7: the \textbf{inverse of the exponential, with the
    domain/range swap}. It covers \textbf{A2.F.2a} (domain, range, zeros, intercepts),
    \textbf{A2.F.2c} (increasing/decreasing intervals), \textbf{A2.F.2d/e} (absolute and relative
    extrema --- here, the answer is \emph{none}, and students must justify it), \textbf{A2.F.2f}
    (evaluate $f(x)$ and interpret it in context --- the doubling model), \textbf{A2.F.2g} (end
    behavior), and \textbf{A2.F.2h} (determine the equation of a vertical asymptote, which names
    logarithmic explicitly). \textbf{A2.F.2b} (compare and contrast the characteristics of function
    families) drives the exponential-vs.-logarithmic table and the Activity Tier~A contrast.
    \textbf{A2.F.1a} (distinguish between the graphs of the parent functions) is the lesson's
    backbone, and \textbf{A2.F.1e} (compare and contrast graphs, tables, and equations) is assessed
    in the Tier~A table and the Homework extension. The \emph{inverse} strand covers
    \textbf{A2.F.2j} (graph the inverse as a reflection over $y=x$) and previews \textbf{A2.F.2i}
    without yet requiring the algebraic method.
\end{tcolorbox}

\renewcommand{\arraystretch}{1.4}
\begin{skillbox}[Priority Ideas \& Skills]{goldbox}
\begin{tabularx}{\linewidth}{@{}X|X@{}}
\textbf{Priority Ideas \& Skills} & \textbf{Key Understandings (the ``why'')} \\[2pt]
\hline
\vspace{-6pt}
\begin{itemize}[leftmargin=1.1em, nosep, after=\vspace{2pt}]
  \item Read $\log_b x$ as \textbf{the exponent you put on $b$ to get $x$}, and use it to evaluate a point on a graph.
  \item Find the \textbf{domain} by solving \textbf{argument} $>0$ --- including the sign-flip case that opens the graph to the \textbf{left}.
  \item Write the equation of the \textbf{vertical asymptote}, and explain why it is never included in the domain.
  \item Read \textbf{range}, \textbf{intercepts}/\textbf{zeros}, \textbf{increasing/decreasing} and \textbf{positive/negative} intervals off both cases, $b>1$ and $0<b<1$.
  \item State \textbf{end behavior} in correct notation --- $x\to+\infty$ on one side, $x\to 0^{+}$ on the other.
  \item Justify that a logarithmic graph has \textbf{no extrema} at all.
  \item Explain the \textbf{inverse relationship} with $y=b^x$: reflection over $y=x$, domain/range swap, horizontal asymptote $\to$ vertical asymptote.
  \item \textbf{Compare and contrast} logarithmic with exponential, and both with the earlier families.
\end{itemize}
&
\vspace{-6pt}
\begin{itemize}[leftmargin=1.1em, nosep, after=\vspace{2pt}]
  \item \emph{The argument runs the show.} A power of a positive base is always positive, so a logarithm has nothing to return for a non-positive argument. Every domain question traces back to that one fact.
  \item Unit~5 made domain something you \emph{solve for}. Unit~7 keeps that --- but changes $\ge$ to $>$, and that single character changes an endpoint into an asymptote.
  \item A square root got to \emph{stand on} its edge. A logarithm's edge is \textbf{missing}: the graph approaches it forever and never arrives.
  \item ``End behavior'' has meant $x\to\pm\infty$ all year. Here one side is $x\to 0^{+}$ --- the input runs down to a \emph{finite} boundary, not out to infinity.
  \item Nothing about this family is arbitrary: it is $y=b^x$ read backwards. Domain, range, asymptote, and intercepts are all Unit~6 facts with the coordinates traded.
  \item $y=x^2$ had to be cut in half before $\sqrt{x}$ could undo it. $y=b^x$ never repeats an output, so \emph{all} of it reflects --- no restriction required.
\end{itemize}
\end{tabularx}
\end{skillbox}

\begin{skillbox}[{Vocabulary, Concepts, \& Theorems}]{sky}
\renewcommand{\arraystretch}{1.3}
\begin{tabularx}{\linewidth}{@{}l X@{}}
\textbf{Logarithmic function} & A function $f(x)=\log_b x$ that returns \emph{the exponent you must put on $b$ to get $x$}; the inverse of $f(x)=b^x$. \\
\textbf{Base $b$} & The number being raised to a power, written low after ``$\log$''. As in Unit~6, $b>0$ and $b\ne1$. \\
\textbf{Argument} & The expression inside the logarithm --- the value you are trying to reach with a power of $b$. \\
\textbf{Domain rule} & The argument must be \emph{strictly} positive, so the domain is the solution set of \emph{argument} $>0$. Note the strict inequality: unlike Unit~5, the boundary is excluded. \\
\textbf{Vertical asymptote} & The vertical line at the edge of the domain, found by setting \emph{argument} $=0$. The graph approaches it but never touches it. Revisits the Unit~4 idea in a new family. \\
\textbf{Common logarithm} & $\log x$ with no base written means base $10$. (Used in notation only today; developed in 7.1.) \\
\textbf{Logarithmic parent, $b>1$} & $y=\log_2 x$: domain $(0,\infty)$, range all reals, VA $x=0$, $x$-intercept $(1,0)$, no $y$-intercept, increasing, no extrema. \\
\textbf{Logarithmic parent, $0<b<1$} & $y=\log_{1/2} x$: identical domain, range, asymptote, and intercepts --- but \emph{decreasing}, so the end behavior reverses. \\
\textbf{Inverse relationship} & Two functions that undo each other; inputs and outputs trade places, graphs reflect over $y=x$, and domain and range swap. A horizontal asymptote becomes a vertical one. \\
\textbf{Why no restriction is needed} & $y=b^x$ never sends two inputs to the same output, so all of it can be reflected --- unlike $y=x^2$, which Unit~5 had to cut to $x\ge0$ first. \\
\end{tabularx}
\end{skillbox}

\begin{fixedskillbox}[Activate Prior Knowledge \& Spiral Review]{forestbg}
    The Warm-Up rehearses the three skills this lesson stands on, one per new idea: \textbf{(1)}
    evaluating $2^3$, $2^0$, $2^{-2}$ and then running the same machine \emph{backwards}
    ($2^{\square}=32$, $2^{\square}=\frac18$, $2^{\square}=-8$) --- so students compute two logarithms
    before the word exists, and discover for themselves that exactly one of the three backwards
    questions is impossible; \textbf{(2)} reading domain, range, horizontal asymptote, and both
    intercepts off the Unit~6 parent $y=2^x$ --- every one of these gets \emph{reflected} into a new
    fact in Section~4 of the notes; and \textbf{(3)} completing a table for $y=2^x$ and then the same
    table with the rows traded, closing with ``look only at the inputs of the second table --- what do
    they all have in common?'' (they are all positive), which hands students the domain $x>0$ before
    they see a single logarithmic graph. Debrief all three aloud, but do \emph{not} supply the words
    ``logarithm,'' ``argument,'' or ``inverse'' --- let the Hook graphs name them.
\end{fixedskillbox}

\begin{skillbox}[Hook]{forestbg}
    Put two graphs on the board side by side: the familiar $y=2^x$ and the brand-new $y=\log_2 x$.
    Open with the observation that frames the whole unit: \textbf{``These are the same picture. One of
    them has been tipped over. Which line would you fold the page along to turn one into the
    other?''} Let students hunt --- most will point at the diagonal before they can name $y=x$. Then
    press on the asymptote: \textbf{``The exponential graph never touched the $x$-axis. This new one
    never touches the $y$-axis. What happened to that flat line when the page was folded?''} (It stood
    up. The horizontal asymptote became a vertical one.) Then pull the domain out of the Warm-Up:
    \textbf{``Item 1 already told you why the new graph has nothing on the left. Why was
    $2^{\square}=-8$ impossible?''} (A power of a positive base is always positive --- so there is no
    exponent to report for a negative input.) Land the point: \emph{Unit~6 asked ``I have the exponent
    --- what is the value?'' This unit asks it backwards: ``I have the value --- what was the
    exponent?'' Today we learn to read the graph that answers that question. Naming it comes first;
    computing with it is Lesson 7.1.} Say the staging out loud --- it is the same move that let
    students read $\sqrt{x}$ in 5.0 before meeting rational exponents in 5.1.
\end{skillbox}

\begin{skillbox}[Lesson]{forestbg}
\begin{multicols}{2}
\textbf{1. Base, argument, and where the graph is allowed to exist.} Give the definition as a
sentence, not a formula: $\log_b x$ is \emph{the exponent you put on $b$ to get $x$}. Name the two
parts (\textbf{base} $b$, written low, with $b>0$ and $b\ne1$ inherited from Unit~6; \textbf{argument},
the expression inside). Read six values off the definition with no calculator --- $\log_2 8$,
$\log_2 1$, $\log_3 9$, $\log_2\frac14$, $\log_5 5$, and then $\log_2(-8)$, which stops the room.
Connect it to Warm-Up 1(f) and state the rule the unit rests on: \emph{set the argument $>0$ and
solve}. Work $\log_2(x-1)$, $\log_3(x+4)$, and $\log(5-x)$ (flip! --- and flag that this one runs
\emph{left}). Then spend a full minute on the character that changed: Unit~5 allowed radicand $\ge0$,
so a square root \emph{stood on} its edge; a logarithm demands argument $>0$ \emph{strictly}, so the
edge is missing and the boundary is a \textbf{vertical asymptote}. Round brackets, every time.

\textbf{2. The parent with $b>1$.} Build $y=\log_2 x$ from the table
$(\frac14,-2),(\frac12,-1),(1,0),(2,1),(4,2),(8,3)$ --- students can fill it entirely by asking ``2 to
what power?'' Read every feature: domain $(0,\infty)$, range \emph{all reals}, VA $x=0$, $x$-intercept
$(1,0)$, \emph{no} $y$-intercept, increasing on the whole domain, negative on $(0,1)$ and positive on
$(1,\infty)$, and \textbf{no extrema at all}. Spend real time on end behavior: as $x\to+\infty$,
$y\to+\infty$ --- but so slowly that $x$ must \emph{double} to lift $y$ by one. And on the other side,
insist on the notation: the input does \emph{not} run to $-\infty$; it runs down to the asymptote, so
as $x\to 0^{+}$, $y\to-\infty$. Write $0^{+}$ on the board and make students copy it.

\columnbreak

\textbf{3. The case $0<b<1$.} Same treatment for $y=\log_{1/2}x$, and let students discover how
little changes: domain, range, asymptote, both intercepts, and the absence of extrema are
\emph{identical}. Only the direction flips --- decreasing, with the end behavior and the
positive/negative intervals reversing with it. Close the section with the question that is really an
A2.F.2b item: \emph{which four rows are the same, and what single fact do they all come from?}
(Whatever the base, $b^{\text{anything}}$ is positive and $b^0=1$.) This mirrors 6.0's growth-vs.-decay
pair exactly.

\textbf{4. Where these parents come from --- the inverse relationship.} Return to Warm-Up item~3: the
two tables were the same table with the rows traded. Name it: raising to a power and asking ``to what
power?'' \textbf{undo} each other, and on a graph undoing is a \textbf{reflection over $y=x$}. Draw
$y=2^x$ and $y=\log_2 x$ on one grid with the diagonal. Then get all four payoffs: the domain of
$\log_2 x$ is $(0,\infty)$ \emph{because that was the range of} $2^x$; the range is all reals
\emph{because that was the domain}; the horizontal asymptote $y=0$ reflects into the vertical
asymptote $x=0$; and $(0,1)$ with no $x$-intercept reflects into $(1,0)$ with no $y$-intercept. Finish
with the Unit~5 contrast: $y=x^2$ needed a restricted domain before $\sqrt{x}$ could undo it, but an
exponential never repeats an output, so all of it reflects.

\textbf{5. The full feature list, and the contrast.} Read every characteristic off the anchor
$f(x)=\log_2(x-1)$ (domain $(1,\infty)$, range all reals, VA $x=1$, zero $(2,0)$, no $y$-intercept,
increasing throughout, negative on $(1,2)$, no extrema, $x\to1^{+}\Rightarrow f\to-\infty$). Finish
with the \textbf{A2.F.2b} table comparing logarithmic with exponential row by row --- every row is the
row above it read backwards --- and one sentence against the earlier units: a polynomial was defined
everywhere; a rational function lost scattered inputs; a radical function lost half the line but kept
its edge; a logarithmic function loses half the line \emph{and} the edge. Preview Lesson~7.1,
\emph{Introduction to Logarithms}.
\end{multicols}
\end{skillbox}

\begin{skillbox}[Active Monitoring]{redbox}
Circulate during the notes practice and Tier work. Watch for and correct:
\begin{itemize}[leftmargin=1.2em, nosep]
  \item writing a log domain with a \emph{square} bracket ($[1,\infty)$ for $\log_2(x-1)$) --- the Unit~5 habit carried over; this is the signature error of the lesson, and the bracket is the tell that the student thinks the boundary is an endpoint rather than an asymptote;
  \item answering ``as $x\to-\infty$'' for the left end --- there is no left end; the input runs down to $0^{+}$, a \emph{finite} boundary, and the notation must say so;
  \item solving $5-x>0$ as $x>5$ --- forgetting the flip, and therefore claiming the graph runs the wrong way (send them back to the Unit~1 inequality rule);
  \item hunting for a $y$-intercept on a parent log graph and being unwilling to write ``none'' --- then over-correcting in Tier~A, where $\log_2(x+2)$ genuinely \emph{has} one;
  \item confusing the asymptote with the $x$-intercept: both sit near the left of the picture, and $(1,0)$ is easy to mislabel as ``$x=1$ is the asymptote'';
  \item hunting for a maximum or minimum out of reflex --- a logarithmic graph has neither, and students should be able to say \emph{why} (it falls forever near the asymptote and climbs forever to the right);
  \item writing an interval that leaks outside the domain (``increasing on $(-\infty,\infty)$'' for $\log_2(x-1)$).
\end{itemize}
Cold-call prompts: ``What is the argument, and what did you set it greater than?'' ``Did you write a
round bracket or a square one --- and which one is right here, and why?'' ``What is the asymptote's
equation, and how did you get it from the equation alone?'' ``What happens on the left side of this
graph? Say it exactly, in notation.'' ``Where did the domain $(0,\infty)$ come from --- what was it in
Unit~6?''
\end{skillbox}

\begin{skillbox}[Group Work \& Differentiation]{redbox}
\textbf{Reading Logarithmic Graphs} activity (tiered; mirrors the notes).
\begin{multicols}{3}
\textbf{Tier R --- Remediate.} From one pre-drawn graph of $g(x)=\log_3 x$: identify the argument and
solve for the domain, name the vertical asymptote and explain why the graph can never reach it, give
the range and $x$-intercept, then \emph{attempt} the $y$-intercept (``3 to what power gives 0?'') and
conclude there is \textbf{none}. Closes on the positive/negative intervals and on end behavior in
both directions --- with the left one required to be written as $x\to 0^{+}$. Send here any student
who wrote square brackets in the notes.

\columnbreak
\textbf{Tier A --- Approaching.} A nine-row feature table for two functions chosen so that most rows
differ: $p(x)=\log_2(x+2)$ (VA $x=-2$, increasing, and it \emph{does} have a $y$-intercept, $(0,1)$)
beside $q(x)=\log_{1/2}x$ (VA $x=0$, decreasing, no $y$-intercept). Then two \emph{justifications}:
point to the feature of each equation that decides whether a $y$-intercept exists, and explain how to
find each asymptote from the equation alone. Expect the productive argument that a log graph can
never have a $y$-intercept --- $p$ is the counterexample; pause the room for it.

\columnbreak
\textbf{Tier E --- Extension.} Equation-only work: a four-row table (argument condition / domain /
asymptote) over $\log_2(x-5)$, $\log_3(x+4)$, $\log(2x-6)$, and $\log_5(9-x)$, plus ``which one runs
left, and how do you know before graphing?'' Then a doubling model $t(n)=\log_2 n$ (hours elapsed
when a colony is $n$ times its original size): domain from algebra \emph{and} context, the meaning of
the $x$-intercept $(1,0)$, why multiplying the size by 16 adds only 4 hours, rewriting
$t=\log_2 n$ as $n=2^t$ and recognizing it as the \textbf{inverse} --- then explaining what the
asymptote at $n=0$ \emph{means} about the bacteria, which is the strongest item on the page.
\end{multicols}
\end{skillbox}

\begin{skillbox}[Individual Work \& Assessment]{redbox}
\textbf{Exit Ticket} (independent, no notes): from one graphed logarithmic function
$f(x)=\log_2(x+4)$, give the domain and range in interval notation, the vertical asymptote, both
intercepts (it \emph{has} a $y$-intercept, $(0,2)$ --- deliberately, so ``no $y$-intercept'' cannot be
answered by reflex), the increasing interval, the absence of extrema, and end behavior in \emph{both}
directions with correct $x\to-4^{+}$ notation; then \textbf{explain} why $f$ has a restricted domain
while $y=2^x$ does not, \emph{and} why that boundary is an asymptote rather than a Unit~5 endpoint;
then an \textbf{SOL-style multiple-choice} item on the domain of $h(x)=\log_5(3-x)$ (answer $x<3$)
whose distractors each break exactly one idea (forget the inequality flip / use the Unit~5 non-strict
bracket / treat a logarithm's argument as unrestricted). Collect and sort into ``got it /
sign-flip slip / asymptote-as-endpoint / domain confusion'' to decide whether Lesson~7.1 opens with a
re-teach --- the third pile is the biggest and the cheapest to fix, but the fourth matters most,
since 7.1 re-derives the domain algebraically from the definition of a logarithm.
\end{skillbox}

\begin{skillbox}[Reinforcement \& Extension]{goldbox}
\textbf{Homework:} a four-row equation-only table (argument condition / domain / vertical asymptote /
which way it runs) over $\log_3(x-2)$, $\log(x+7)$, $\log_2(4x)$, and $\log_5(10-2x)$; a full feature
read off $m(x)=\log_2 x-3$, chosen because the vertical shift moves the $x$-intercept to $(8,0)$ but
leaves the asymptote at $x=0$ --- and the closing question asks students to say why; a matched full
feature read off the decreasing $n(x)=\log_{1/3}x$, closing with a list of every feature it shares
with $\log_3 x$; and one justification item on why every parent log graph passes through $(1,0)$ and
never crosses the $y$-axis, using the words \emph{exponent} and \emph{asymptote}.
\textbf{Extension:} build a logarithmic graph from an exponential table --- given the table for
$y=3^x$, trade the rows, read the domain and range off the result, reflect the horizontal asymptote
$y=0$ over $y=x$ and name what it becomes, predict both intercept facts from the coordinate swap, and
finally explain what is true about an exponential graph that let \emph{all} of it be reflected when
$y=x^2$ had to be cut in half first. \textbf{Preview:} Lesson~7.1, \emph{Introduction to Logarithms}
--- $\log_b x=y$ and $b^y=x$ as one statement read two ways, evaluating by asking ``$b$ to what
power?'', the common log, the four identities, and change of base for calculator work, which is also
where today's domain rule gets its full algebraic argument.
\end{skillbox}
\end{document}
